\chapter{Introducció}
\label{ch:introduccio}

\section{Context}
\label{sec:context}

La quantitat d'informació personal que una persona genera i consumeix ha
crescut de manera exponencial durant l'última dècada: correus, notes,
documents, captures, articles desats per llegir més tard i converses
esteses. Aquesta informació viu fragmentada en múltiples eines
(Notion, Obsidian, Google Keep, Evernote, Apple Notes), sense un punt
d'accés unificat ni una interfície conversacional que permeti
interrogar-la amb preguntes en llenguatge natural.

La irrupció dels grans models de llenguatge (LLM) durant el 2022--2025
ha habilitat un patró d'interacció nou: el xat contextual basat en
\emph{Retrieval-Augmented Generation} (RAG), que permet fer preguntes
sobre un corpus privat indexat semànticament. Aquest patró ha passat de
curiositat acadèmica a producte de producció: al 2026, eines com
Notion~AI, Mem, Obsidian Copilot o els \emph{Projects} de ChatGPT ho
ofereixen comercialment. La diferència entre un prototip i un producte
defensable és, però, la seva \textbf{arquitectura multi-tenant segura}
(cada usuari veu només les seves dades) i la seva
\textbf{interoperabilitat amb agents d'IA} (altres eines poden operar
sobre les dades amb el consentiment de l'usuari).

L'any 2024, Anthropic va publicar el \textbf{Model Context Protocol}
(MCP) \autocite{anthropic2024mcp}, un protocol obert per estandarditzar
com els hosts de LLM (Claude Desktop, Cursor, agents propis)
\emph{descobreixen} i \emph{criden} eines, recursos i \emph{prompts}
externs. Al llarg del 2025 i 2026, MCP s'ha consolidat com la capa
d'integració dominant per a la IA amb agents, amb implementacions de
referència d'AWS \autocite{aws2026mcp}, Cloudflare
\autocite{cloudflare2026agents} i OWASP \autocite{owasp2026mcp}, i amb
una especificació de seguretat pròpia
\autocite{mcpspec2026security,mcpspec2026auth} que fixa
OAuth 2.1 com a esquema d'autorització.

\textbf{Synapse Notes} és el projecte desenvolupat com a Treball de
Fi de Grau al llarg del curs 2025--2026. L'execució s'ha estructurat
en dues parts complementàries, cadascuna amb els seus propis
entregables i criteris de verificació:

\begin{itemize}
  \item \textbf{Part~A. Plataforma base} (octubre~2025 -- abril~2026).
    Disseny i implementació d'un SaaS multi-tenant de notes amb xat
    contextual RAG. \emph{Stack:} Next.js~15 amb App Router, React~19,
    TypeScript, Tailwind~4 i Shadcn/UI; Supabase com a capa de dades
    (PostgreSQL amb l'extensió pgvector) i d'autenticació; Google
    Gemini \texttt{embedding-001} per als embeddings i Vercel AI SDK
    per al xat en \emph{streaming} amb \emph{tool calling}. La
    plataforma implementa OAuth de Google i GitHub, CRUD de notes,
    internacionalització (CA/ES/EN) i desplegament continu a Vercel.
    Tot l'accés a dades queda protegit per polítiques Row-Level
    Security (RLS) de PostgreSQL.

  \item \textbf{Part~B. Extensió crítica per a agents d'IA}
    (abril -- juny~2026). Aportació de recerca i seguretat del
    treball: servidor \textbf{MCP multi-tenant} amb OAuth~2.1 sobre
    Supabase Auth, 3~agents en segon pla sobre Edge Functions de
    Supabase amb \texttt{pg\_cron}, i un model d'amenaces formal
    basat en la \emph{Lethal Trifecta}
    \autocite{willison2025trifecta} amb avaluació empírica de les
    mitigacions. Des d'abril de 2026, el model de xat passa a
    Claude Haiku~4.5 per homogeneïtzar proveïdor amb la passada de
    seguretat del servidor MCP.
\end{itemize}

\section{Motivació}
\label{sec:motivacio}

La Part~A del TFG consolida una plataforma RAG funcional i els patrons
de SaaS multi-tenant (autenticació OAuth, aïllament per RLS, cerca
vectorial i xat en \emph{streaming} amb eines). La Part~B construeix
sobre aquesta base l'aportació que fa el treball defensable davant el
tribunal: \textbf{convertir la plataforma en una superfície accessible
per agents d'IA de manera segura}. La pregunta de recerca que guia la
Part~B i vertebra la memòria és:

\begin{quote}
  \emph{Com es dissenya, implementa i avalua empíricament un servidor MCP
  multi-tenant sobre una base de coneixement personal, de manera que
  agents autònoms puguin operar-hi sense violar l'aïllament entre
  tenants ni caure en atacs d'injecció indirecta de prompt?}
\end{quote}

La motivació és tripla:

\begin{enumerate}
  \item \textbf{Oportunitat tècnica.} MCP és un estàndard de
    \emph{2026}: no estava disponible quan es va dissenyar Synapse Notes,
    i la literatura sobre servidors MCP \emph{multi-tenant} és molt
    recent \autocite{aws2026mcp,gitguardian2026oauth}. Aplicar-lo sobre
    una base de producte permet contrastar el patró emergent amb les
    pràctiques establertes de SaaS multi-tenant amb RLS
    \autocite{supabase2026rag,tigerdata2026multitenant,nile2026multitenant}.

  \item \textbf{Rellevància en seguretat.} Durant el gener de 2026 es
    van divulgar públicament exploits de la denominada ``Lethal
    Trifecta'' \autocite{willison2025trifecta} en eines AI
    industrialitzades (IBM Bob, Superhuman AI, Notion AI i Claude
    Cowork d'Anthropic), totes elles en una sola setmana. La trifecta
    estableix que un agent amb accés a \emph{dades privades},
    \emph{contingut no confiable} i \emph{comunicació externa} és
    incondicionalment vulnerable a la injecció indirecta de prompt. El
    TFG aplica aquest marc eina a eina al servidor MCP resultant i
    mesura empíricament l'efecte de les mitigacions amb una suite
    automatitzada de red-team \autocite{promptfoo2026trifecta}.

  \item \textbf{Alineació formativa.} Els objectius del TFG cobreixen
    competències centrals del Grau en Enginyeria Informàtica de la
    URV: \emph{Seguretat de Xarxes} (OAuth~2.1, JWT, modelatge
    d'amenaces), \emph{Sistemes Distribuïts} (arquitectura
    client/servidor MCP, agents en segon pla amb \texttt{pg\_cron}),
    \emph{Bases de Dades Avançades} (RLS, pgvector, HNSW),
    \emph{Intel·ligència Artificial} (RAG, LLMs, \emph{tool calling}),
    \emph{Desenvolupament Avançat d'Aplicacions Web} (Next.js, React,
    App Router) i \emph{Anàlisi i Disseny d'Aplicacions} (UML, casos
    d'ús, requisits). Aquesta triangulació justifica el TFG a l'ETSE.
\end{enumerate}

\section{Ús esperat i usuaris objectiu}
\label{sec:usuaris}

L'usuari final de Synapse Notes és una persona amb una base de
coneixement personal no trivial que vol interrogar-la amb un agent
d'IA. Els tres fluxos d'ús esperats són:

\begin{itemize}
  \item \textbf{Xat RAG integrat} (Part~A). L'usuari accedeix a la UI
    web i pregunta sobre les seves notes amb llenguatge natural. El
    sistema recupera els fragments rellevants via pgvector i els passa
    al LLM (Claude Haiku~4.5) en \emph{streaming}, amb \emph{tool
    calling} disponible (\texttt{getNotesByTag}).

  \item \textbf{Agents externs via MCP} (Part~B). L'usuari connecta
    el seu client MCP preferit (Claude Desktop, Cursor, un agent propi)
    a la seva instància de Synapse. El client autentica via OAuth~2.1,
    rep un JWT per tenant i invoca eines com \texttt{search\_notes} o
    \texttt{summarise\_notes}.

  \item \textbf{Agents en segon pla} (Part~B). Tres agents
    (\texttt{embedding-backfill}, \texttt{auto-tag},
    \texttt{weekly-digest}) operen sense intervenció de l'usuari, sota
    \texttt{pg\_cron} a Supabase, i escriuen els seus resultats a
    taules auditables
    \autocite{supabase2026cron,supabase2026largejobs}.
\end{itemize}

\section{Abast del treball}
\label{sec:abast}

L'abast del TFG és el conjunt del projecte Synapse Notes, organitzat en
les dues parts presentades a la secció~\ref{sec:context}. Les
aportacions concretes s'enumeren a continuació.

\textbf{Part~A. Plataforma base} (desenvolupament continu d'octubre~2025
a abril~2026):

\begin{enumerate}
  \item Autenticació OAuth multi-proveïdor (Google i GitHub) sobre
    Supabase Auth, amb sessió persistent i protecció de rutes a
    Next.js.
  \item Model de dades multi-tenant amb polítiques RLS sobre
    \texttt{notes}, \texttt{chats} i \texttt{messages}, i RPC
    \texttt{match\_notes} per a cerca vectorial amb pgvector.
  \item CRUD de notes i xat en \emph{streaming} via Vercel AI SDK amb
    \emph{tool calling} (\texttt{getNotesByTag}), integrat amb el
    context RAG.
  \item UI internacionalitzada (CA/ES/EN) amb Shadcn/UI, paleta
    d'ordres, mode fosc i desplegament continu a Vercel.
\end{enumerate}

\textbf{Part~B. Extensió crítica per a agents d'IA} (desenvolupament
intensiu del 19~d'abril al 5~de juny de 2026):

\begin{enumerate}
  \item Un servidor \textbf{MCP multi-tenant} amb 6 eines
    (\texttt{search\_notes}, \texttt{get\_note}, \texttt{create\_note},
    \texttt{update\_note}, \texttt{tag\_notes},
    \texttt{summarise\_notes}), transport Streamable HTTP i
    autenticació OAuth 2.1 sobre Supabase Auth.

  \item \textbf{3 agents en segon pla} desplegats com a Edge Functions
    de Supabase amb \texttt{pg\_cron}: \texttt{embedding-backfill}
    (cada 15~min), \texttt{auto-tag} (horària) i \texttt{weekly-digest}
    (dominical).

  \item Un \textbf{model d'amenaces formal} basat en la Lethal
    Trifecta, aplicat eina a eina, amb mitigacions implementades
    (etiquetatge de procedència, segregació de capacitats,
    aprovacions humanes, filtre de sortida) i avaluades empíricament
    amb una suite Promptfoo de $\geq 15$ casos red-team i una suite
    d'aïllament RLS de $\geq 15$ casos.
\end{enumerate}

Les exclusions explícites (aplicació mòbil, workspaces d'equip,
facturació, panell d'administració, cerca full-text, fine-tuning de LLM,
Supabase auto-allotjat) queden justificades com a treball futur a la
valoració final. L'abast de la Part~B es \emph{congela} al final de la
primera setmana com a mesura de contenció d'\emph{scope creep}, un dels
dos riscos alts identificats al capítol~\ref{ch:planificacio}.

\section{Organització del document}
\label{sec:organitzacio}

La memòria segueix l'estàndard ADA de la URV, amb 17 seccions:

\begin{itemize}
  \item \textbf{Capítols 1--3} (aquesta introducció, paraules clau,
    objectius). Fixen el context, el vocabulari i els criteris de
    verificació.
  \item \textbf{Capítol~\ref{ch:planificacio}} (planificació).
    Fases de Part~A i Part~B, calendari intensiu de 47 dies per a la
    Part~B, Gantt i taula de riscos.
  \item \textbf{Capítol~\ref{ch:requisits}} (requisits). 19 casos
    d'ús funcionals (4 de Part~A + 15 de Part~B) i 11 requisits
    no funcionals.
  \item \textbf{Capítol~8 (disseny)}. Arquitectura C4, disseny de BD
    amb polítiques RLS, disseny del servidor MCP i \emph{disseny de
    seguretat} (contribució principal, Lethal Trifecta eina a eina).
  \item \textbf{Capítol~9 (implementació)}. Justificació tecnològica,
    fragments de patrons clau i problemes reals trobats.
  \item \textbf{Capítol~\ref{ch:avaluacio}} (avaluació). Tests
    funcionals, aïllament RLS, red-team Promptfoo, benchmarks i cas
    d'estudi amb Claude Desktop.
  \item \textbf{Capítols~11--13} (costos, legislació, ètica). Costos
    de personal i cloud, anàlisi RGPD/LOPDGDD/LSSI i implicacions
    ètiques.
  \item \textbf{Capítols~14 i annexos.} Valoració personal,
    bibliografia APA~7 i annexos tècnics (instal·lació, ús, schemes
    d'eina, YAML de Promptfoo).
\end{itemize}
